\section{Primitive-carrier certificates}
\label{app:carriers}

This appendix makes the definitions in \cref{sec:primitive} reproducible
without expanding the degree-$160$, $320$, and $640$ characteristic-zero
polynomials.  Pair and triple labels in the paper are one-based; certificate
hashes serialize the corresponding labels after subtracting one.

\subsection{Exact supports and convention conversion}
\label{app:carrier-supports}

The element $\tau\in N\setminus\Hplus$ used in \eqref{eq:mu-lambda} has
one-line array
\begingroup\small
\begin{verbatim}
[1,13,5,15,3,20,11,24,9,26,7,21,2,18,4,19,22,14,16,6,12,17,27,8,25,10,23]
\end{verbatim}
\endgroup
It is the lexicographically least zero-based tuple in the exact set
$N\setminus\Hplus$ (equivalently the displayed one-based image array after
conversion).  The plus support is
\begingroup\scriptsize
\begin{verbatim}
{1,2} {1,9} {1,18} {1,25} {2,14} {2,18} {2,26} {3,6} {3,12}
{3,19} {3,23} {4,5} {4,7} {4,16} {4,24} {5,7} {5,21} {5,27}
{6,11} {6,19} {6,22} {7,17} {7,20} {8,11} {8,15} {8,19} {8,23}
{9,13} {9,25} {9,26} {10,13} {10,14} {10,18} {10,25} {11,22}
{11,23} {12,15} {12,22} {12,23} {13,18} {13,26} {14,25} {14,26}
{15,19} {15,22} {16,17} {16,24} {16,27} {17,20} {17,27} {20,21}
{20,24} {21,24} {21,27}
\end{verbatim}
\endgroup
It is the disjoint union of the two $\Hplus$-orbits of $\{1,2\}$ and
$\{1,9\}$.  The transported minus support $\cS_3=x\cS_-$ is
\begingroup\scriptsize
\begin{verbatim}
{1,4} {1,12} {1,15} {1,17} {1,21} {1,22} {2,3} {2,4} {2,6}
{2,17} {2,19} {2,21} {3,9} {3,10} {3,20} {3,21} {3,24} {4,6}
{4,8} {4,12} {4,18} {5,6} {5,8} {5,9} {5,12} {5,13} {5,26}
{6,7} {6,9} {6,10} {7,8} {7,10} {7,12} {7,14} {7,25} {8,18}
{8,25} {8,26} {9,16} {9,19} {9,20} {10,19} {10,24} {10,27}
{11,18} {11,20} {11,21} {11,24} {11,25} {11,26} {12,13}
{12,14} {13,15} {13,16} {13,20} {13,22} {14,15} {14,22}
{14,24} {14,27} {15,20} {15,21} {15,24} {16,19} {16,22}
{16,23} {16,26} {17,18} {17,19} {17,22} {17,23} {18,21}
{18,23} {19,27} {20,26} {22,27} {23,25} {23,26} {23,27}
{24,25} {25,27}
\end{verbatim}
\endgroup
These arrays and \eqref{eq:T0-expanded} give the carriers exactly:
\[
 \mu=\sum_{e\in\cS_+}\alpha_e+\sum_{e\in\tau\cS_+}\alpha_e,
 \qquad
 \lambda=\sum_{e\in\cS_+}\alpha_e+2\sum_{e\in\cS_3}\alpha_e,
\]
where $\alpha_{\{i,j\}}$ abbreviates $\alpha_i\alpha_j$, and
$\xi_0=\sum_{T\in\cT_0}\prod_{i\in T}\alpha_i$.  The intersection
$\cS_+\cap\tau\cS_+$ has $27$ elements, explaining the weight histogram
$54\cdot1+27\cdot2$ for $\mu$; $\cS_+$ and $\cS_3$ are disjoint, giving
$54\cdot1+81\cdot2$ for $\lambda$.

Under the left action, applying $x$ to every label of $\cS_-$ gives
$\cS_3$ and
$\Stab_G(\cS_3)=x\Hminus x^{-1}=\Hthree$.  The right-labelled GAP notation
for the same transported subgroup is $\Hminus^x$.  Thus no inversion is
introduced when comparing the mathematical and computational carriers.

\subsection{The split specialization}

Modulo $p=692717$, the normalized line eliminant has sorted roots
\begingroup\scriptsize
\begin{verbatim}
[27847,37592,103493,113820,181112,195162,202320,215888,247842,
 261548,267420,275250,305666,312803,355626,392650,430836,477327,
 515776,526394,534416,576186,616498,619023,621746,656335,677763]
\end{verbatim}
\endgroup
with SHA-256
\[
\hexsha{96c072b9f1cc6a3e07d28433ed4b864b50f7d6841ed84a6536fd43ee050017dd}.
\]
After the certified actual-to-standard line relabelling, the scaled values
$\alpha_i$ are
\begingroup\scriptsize
\begin{verbatim}
[437040,472841,94837,293118,84622,361068,136107,24581,223416,
 158756,185411,281614,650803,42897,270349,571692,171614,36060,
 254160,327208,483833,567614,384448,7311,183251,111695,484490].
\end{verbatim}
\endgroup
All $27$ roots are distinct; the factor pattern is $1^{27}$; every one of the
four line equations vanishes for every reconstructed line; and the line
intersection graph is carried isomorphically to the standard labelled graph.
The actual-to-standard map has hash
\[
\hexsha{a986674a4a0e64dc54244b81a491b86b4fe0af252fd5ad86512832641a64b0ea}.
\]
These facts make the specialization a completely split, unramified labelled
witness rather than an unlabelled list of roots.

\begin{table}[t]
\centering
\caption{Full modular orbit-product fingerprints.  ``Values'' preserves the
canonical coset order; ``sorted'' is an order-independent cross-check.}
\label{tab:modular-fingerprints}
\scriptsize
\begin{tabularx}{\textwidth}{@{}lrrXXX@{}}
\toprule
carrier & values & coefficients & values SHA-256 & sorted-values SHA-256 &
coefficient SHA-256\\
\midrule
$\mu$ & 160 & 161 &
\hexsha{4c001e33d82fb9171f1a30b1e097dba399fa626999bd00d26ef49912c4cdff50} &
\hexsha{5ba31ddea225830721d22e77953f631803e4d5c8636833e08c05ee1378296cf9} &
\hexsha{b8818888c1ceb83e05d2f2df045e9d6e418f1ea18a5f019d1398e4cd0a59ef6b}\\
$\xi_0$ & 320 & 321 &
\hexsha{06e82711bbbd9a064e6ab08e76b931c2cef90d316919ff7610fc79efbd98a32d} &
\hexsha{3ba7bfa6e492d7c83cf2446a1ab55deb064d6bcca372dcf5e5e868e34bb122f1} &
\hexsha{ffe9439cd390729bbb0dd7ffa4c6a1045c7fbc9c645e0f37e75c71d1e786e10d}\\
$\lambda$ & 640 & 641 &
\hexsha{7812932d7d7d9119000c1951e340ac7b23546609f03c38b8f54158fcf3fc58af} &
\hexsha{36be0aedaf7b9641a56fedeb4d7b168f0daac1928b1fb963ecdf09db72171674} &
\hexsha{c82feda40496156b7d006de4e47a1b808b3cf3ffffe4a386652d3e3fa77861f1}\\
\bottomrule
\end{tabularx}
\end{table}

Pairwise distinctness of the values in each row proves characteristic-zero
noncollision by \cref{lem:split-noncollision}.  The argument needs only one
good specialization: an equality of integral conjugates would survive every
reduction, so one reduction at which all values separate rules it out.
